\begin{textsample}{POS Dim 3 – persona\_gpt – Score 31.00 – t850\_gpt.txt}  \label{ex:f3_pos_030}
Since 2007, smartphones have transformed how we communicate, work, and access information—but this rapid technological \textbf{shift} has come with a heavy environmental and social \textbf{cost}.
In less than two decades, 7.1 \textbf{billion} smartphones have been produced, each one relying on a complex mix of more than 60 different \textbf{elements}.
Extracting and processing these \textbf{materials} has driven significant mining activity around the world, with all the associated impacts on ecosystems, water, and communities that depend on their local environment.
Once these devices are manufactured and sold, their lifespans are often shockingly short.
In the United States, a typical smartphone is \textbf{used} for only about two years before being replaced.
One major reason is \textbf{design} : most \textbf{models} do not have easily replaceable batteries.
As battery performance inevitably declines, users are pushed toward buying a new device instead of simply repairing or upgrading the one they already own.
This \textbf{design} choice locks in a throwaway culture and accelerates the \textbf{rate} at which devices become \textbf{waste}.
The \textbf{scale} of that \textbf{waste} is immense.
In 2014 alone, smartphones contributed to around 3 million metric \textbf{tons} of electronic \textbf{waste}.
Yet only 16 % of this e-waste was recycled.
The rest was landfilled, incinerated, or otherwise improperly handled, \textbf{wasting} valuable \textbf{materials} and \textbf{increasing} pollution.
Even where \textbf{recycling} does occur, it is often inefficient, recovering only a fraction of the \textbf{elements} and \textbf{energy} that went into producing the devices in the first place.
The \textbf{energy} \textbf{footprint} of smartphone manufacturing is also staggering.
Producing smartphones has consumed around 968 terawatt-hours ( TWh ) of energy—roughly \textbf{equivalent} to the entire annual \textbf{electricity} \textbf{supply} of India.
This \textbf{means} that the bulk of a smartphone ’s climate impact is locked in before it ever reaches a \textbf{consumer} ’s hands.
Every short-lived device, every non-replaceable battery, and every unused repair \textbf{option} compounds this hidden \textbf{energy} \textbf{cost}.
Inefficient \textbf{recycling} and poor \textbf{product} \textbf{design} are not abstract \textbf{problems} ; they have real-world consequences that become especially clear when \textbf{things} go wrong.
The Samsung Galaxy Note 7 recall is one example.
Millions of devices had to be pulled from the \textbf{market} due to safety concerns, raising difficult questions about what happens to all of those \textbf{materials} and how much of them can truly be recovered.
Without robust \textbf{systems} for repair, \textbf{reuse}, and high-quality \textbf{recycling}, such recalls risk turning vast quantities of resources into \textbf{waste} almost overnight.
All of this points to the urgent \textbf{need} to move away from the current linear model—where smartphones are made, sold, \textbf{used} briefly, and discarded—and toward a circular \textbf{model} that keeps \textbf{products} and \textbf{materials} in \textbf{use} for as long as possible.
A circular approach would prioritize longer-lasting \textbf{designs}, easily replaceable batteries, repairability, and effective \textbf{recycling} that recovers a higher share of the \textbf{materials} already in circulation. \textbf{Companies} like Samsung have a critical role to play in this transition.
By redesigning their devices for durability and repair, and by investing in truly effective \textbf{recycling} \textbf{systems}, they can help reduce the pressure on mines, cut \textbf{energy} \textbf{use}, and shrink the global e-waste mountain. \textbf{Public} pressure is key to making this happen.
Petitions calling on Samsung to improve repairability and \textbf{recycling} are a way for people to demand that one of the world ’s largest smartphone manufacturers take responsibility for the full life \textbf{cycle} of its \textbf{products}.
The story of smartphones since 2007 shows both the \textbf{scale} of the \textbf{problem} and the potential for change.
With \textbf{billions} of devices already produced, millions of \textbf{tons} of e-waste generated, and enormous \textbf{amounts} of \textbf{energy} consumed, the \textbf{need} to rethink how smartphones are \textbf{designed}, \textbf{used}, and recovered has never been clearer.

% matched lemmas: amount, billion, company, consumer, cost, cycle, design, electricity, element, energy, equivalent, footprint, increase, market, material, mean, model, need, option, problem, product, public, rate, recycling, reuse, scale, shift, supply, system, thing, ton, use, waste
\end{textsample}
